\documentclass[UTF8,zihao=-4]{ctexart}
\usepackage[a4paper,margin=2.35cm]{geometry}
\usepackage{amsmath,amssymb,bm}
\usepackage{booktabs,longtable,array}
\usepackage{enumitem}
\usepackage{xcolor}
\usepackage{hyperref}
\usepackage{caption}
\usepackage{listings}

\hypersetup{colorlinks=true,linkcolor=blue!55!black,urlcolor=blue!55!black}
\setlist[itemize]{leftmargin=2em,itemsep=0.28em}
\setlist[enumerate]{leftmargin=2em,itemsep=0.28em}
\lstset{basicstyle=\ttfamily\small,breaklines=true,frame=single,columns=fullflexible}

\title{MPC 控制器原理说明\\{\large 从零开始理解 Local MPC、ADMM MPC 与 soft/hard 模式}}
\author{MADRL\_ESS 项目理解笔记}
\date{\today}

\begin{document}
\maketitle
\tableofcontents
\newpage

\section{先不看公式：MPC 到底在做什么}

MPC 的全称是 Model Predictive Control，中文常叫模型预测控制。这个名字里有三个关键词：
\begin{itemize}
  \item Model：我知道系统大致会怎样变化，例如电池充电后 SoC 会升高，放电后 SoC 会降低；
  \item Predictive：我有未来一小段时间的预测，例如未来电价、负荷、PV 发电；
  \item Control：我现在要决定动作，例如电池充不充、EV 充不充、PV 要不要弃光。
\end{itemize}

如果用一句非常直白的话概括：
\begin{quote}
MPC 每一步都先“看一眼未来”，算一份未来计划，但真正执行时只执行计划的第一步。
\end{quote}

为什么只执行第一步？因为预测会变。下一步到来时，系统状态也变了，新的预测也可能不同。MPC 不会死守旧计划，而是重新规划。这种方法叫滚动时域控制。

可以把它想成开车导航。导航会规划未来几十公里，但你真正做的只是当前路口左转还是直行。等你开到下一个路口，导航会根据新路况重新规划。MPC 对电池、EV、PV 做的事情也很像：它规划未来多个时间步，但只把当前一步动作交给环境执行。

\section{优化问题是什么：给完全没学过算法的读者}

MPC 的核心不是神秘的“智能”，而是每一步解一个优化问题。优化问题可以理解为：
\begin{quote}
在一堆允许的选择里，找一个最好的选择。
\end{quote}

一个优化问题通常由三部分组成：
\begin{enumerate}
  \item 决策变量：控制器可以选择的量；
  \item 约束条件：选择时不能违反的规则；
  \item 目标函数：用来判断哪个选择更好的打分函数。
\end{enumerate}

\subsection{决策变量：控制器手里的旋钮}

在本项目中，MPC 可以决定：
\begin{equation}
c_{i,k},\quad d_{i,k},\quad u_{i,k},\quad e_{i,k}.
\end{equation}

它们分别表示：
\begin{itemize}
  \item $c_{i,k}$：第 $i$ 个用户在未来第 $k$ 步给静态电池充多少电；
  \item $d_{i,k}$：第 $i$ 个用户在未来第 $k$ 步让静态电池放多少电；
  \item $u_{i,k}$：第 $i$ 个用户在未来第 $k$ 步丢掉多少 PV 发电；
  \item $e_{i,k}$：第 $i$ 个用户在未来第 $k$ 步给 EV 充多少电。
\end{itemize}

这些量叫决策变量，是因为它们不是已经发生的事实，而是控制器要做出的决定。

\subsection{约束条件：物理世界的边界}

约束条件就是不能违反的规则。例如：
\begin{equation}
0\le e_{i,k}\le P^{ev,\max}_i.
\end{equation}

这句话的意思是：EV 充电功率不能小于 0，也不能超过充电桩和车辆允许的最大功率。约束不是建议，而是物理边界。优化器不能为了省钱或好看而让 EV 充电功率变成负数，也不能让电池能量超过容量。

如果所有约束能同时满足，我们说这个问题是可行的。如果约束互相冲突，例如 EV 只剩 15 分钟就要离家，但要求从 10\% 充到 90\%，最大充电功率又不够，那么问题可能不可行。

\subsection{目标函数：什么叫“最好”}

目标函数是优化器用来比较不同计划的分数。MPC 一般写成：
\begin{equation}
\min J.
\end{equation}

这里 $\min$ 的意思是 minimize，也就是让 $J$ 尽可能小。$J$ 可以是电费、惩罚、设备磨损等多个项目加起来的总成本。

例如，如果电价是 $\tilde p_t$ 欧元/kWh，某一步充电功率是 $P_t$ kW，时间长度是 $\Delta t$ 小时，那么这一小步买电费用近似为：
\begin{equation}
\text{费用}=\tilde p_t P_t\Delta t.
\end{equation}

单位也能对上：
\begin{equation}
\frac{\text{欧元}}{\text{kWh}}\times \text{kW}\times \text{h}
=\text{欧元}.
\end{equation}

所以优化器不是“凭感觉”省钱，而是在每个未来时间步计算成本，然后找总成本最低的一组动作。

\subsection{读 MPC 公式的小技巧}

初学者看到公式容易紧张。其实读 MPC 公式可以按这个顺序：
\begin{enumerate}
  \item 先看变量单位：kW 是功率，kWh 是能量，欧元/kWh 是价格；
  \item 再看等式左右：等式左右应该表示同一类物理量；
  \item 最后看正负号：充电通常增加负荷，放电通常减少负荷，弃光会减少可用 PV。
\end{enumerate}

如果公式能解释成一句物理语言，它通常就不难。例如：
\begin{equation}
E_{k+1}=E_k+\eta P\Delta t
\end{equation}
就是“下一步能量 = 当前能量 + 这一小段时间充进去的能量”。

\section{MPC 和 Rule-based 的区别}

Rule-based 控制器像固定规则表。例如：
\begin{quote}
如果电价低，就充电；如果电价高，就放电。
\end{quote}

这种方法容易理解，但它通常只看当前信息或简单阈值。MPC 不同，它会看未来。例如：
\begin{quote}
现在电价有点低，但两个小时后更低，所以先不充；等两个小时后再充。
\end{quote}

因此，MPC 的核心优势是“有计划”。它不是只问“现在该怎么办”，而是问：
\begin{quote}
如果我现在这样做，接下来几个小时会发生什么？
\end{quote}

\section{本项目中的两种 MPC}

本项目主要有两种 MPC：
\begin{itemize}
  \item Local MPC：每个 agent 独立做自己的优化；
  \item ADMM MPC：每个 agent 还是做自己的优化，但大家通过 ADMM 协调总变压器功率。
\end{itemize}

可以用一个类比理解：
\begin{itemize}
  \item Local MPC 像每户家庭各自省钱，互相不商量；
  \item ADMM MPC 像每户家庭先各自做计划，然后小区协调员检查总功率是否超过变压器限制，超了就让大家改一点。
\end{itemize}

仓库中还存在 Global MISOCP 基线。它像一个中央调度器，直接把全网物理约束放进一个大优化问题。它更完整，但求解更重。本文件重点讲 Local MPC 和 ADMM MPC。

\section{符号说明}

\begin{longtable}{>{\centering\arraybackslash}p{3.0cm}p{3.8cm}p{7.6cm}}
\toprule
符号 & 单位 & 含义 \\
\midrule
\endhead
$i$ & 无 & agent 编号，本项目默认有多个家庭或节点 \\
$k$ & 无 & MPC 预测窗口内的步数，$k=0,\dots,H-1$ \\
$t$ & 无 & 环境当前真实时间步 \\
$H$ & 步 & MPC 向未来看的步数，叫 prediction horizon \\
$\Delta t$ & 小时 & 一个时间步的长度，默认 0.25 小时 \\
$\tilde p_{t+k}$ & 欧元/kWh & 未来第 $k$ 步进口电价 \\
$L_{i,t+k}$ & kW & 第 $i$ 个用户未来第 $k$ 步负荷预测 \\
$G_{i,t+k}$ & kW & 第 $i$ 个用户未来第 $k$ 步 PV 可用功率预测 \\
$c_{i,k}$ & kW & 静态电池充电功率，非负 \\
$d_{i,k}$ & kW & 静态电池放电功率，非负 \\
$P^b_{i,k}$ & kW & 有符号电池功率，$P^b_{i,k}=c_{i,k}-d_{i,k}$ \\
$u_{i,k}$ & kW & PV 弃光功率 \\
$e_{i,k}$ & kW & EV 充电功率 \\
$E^b_{i,k}$ & kWh & 静态电池能量状态 \\
$E^{ev}_{i,k}$ & kWh & EV 电池能量状态 \\
$C^b_i$ & kWh & 静态电池容量 \\
$C^{ev}_i$ & kWh & EV 电池容量 \\
$P^{b,\max}_i$ & kW & 静态电池最大充放电功率 \\
$P^{ev,\max}_i$ & kW & EV 最大充电功率 \\
$\eta_b,\eta_{ev}$ & 无 & 静态电池效率和 EV 充电效率 \\
\bottomrule
\end{longtable}

注意：如果当前真实时间是 $t$，那么预测窗口内的 $k=0$ 就是现在这一小步，$k=1$ 是下一小步，$k=2$ 是再下一小步。

\section{Local MPC：每个用户自己规划}

\subsection{Local MPC 一步一步做什么}

Local MPC 的代码入口是 \texttt{controllers/mpc\_local.py}。每个环境步，它大致做五件事：
\begin{enumerate}
  \item 读取当前状态：电池 SoC、EV SoC、EV 是否在家；
  \item 读取未来预测：未来电价、负荷、PV 可用功率；
  \item 建立优化问题：把可选动作、物理约束和成本函数写进去；
  \item 调用求解器：让求解器找出未来 $H$ 步的低成本计划；
  \item 只执行第一步：把 $k=0$ 的动作交给环境，下一步再重新算。
\end{enumerate}

它的输入可以概括为：
\begin{equation}
\{\tilde p_{t:t+H-1},L_{i,t:t+H-1},G_{i,t:t+H-1}\}.
\end{equation}

它要求解的动作序列是：
\begin{equation}
\{c_{i,k},d_{i,k},u_{i,k},e_{i,k}\}_{k=0}^{H-1}.
\end{equation}

求解完成后，真正送给环境执行的是第一步：
\begin{equation}
P^b_{i,t}=c_{i,0}-d_{i,0},\qquad
G^{curt}_{i,t}=u_{i,0},\qquad
P^{ev}_{i,t}=e_{i,0}.
\end{equation}

这里要特别记住：MPC 虽然算了未来 $H$ 步，但不会一次性全部执行。它只执行当前动作，然后下一步重新预测和优化。

\subsection{为什么电池要拆成充电和放电}

环境里的电池功率可以用一个有正负号的量表示：
\begin{equation}
P^b_{i,k}>0 \text{ 表示充电},\qquad
P^b_{i,k}<0 \text{ 表示放电}.
\end{equation}

但很多优化建模更喜欢非负变量，所以 Local MPC 把电池功率拆成两个变量：
\begin{equation}
P^b_{i,k}=c_{i,k}-d_{i,k}.
\end{equation}

其中：
\begin{equation}
c_{i,k}\ge 0,\qquad d_{i,k}\ge 0.
\end{equation}

如果 $c_{i,k}=5,d_{i,k}=0$，就是充电 5 kW。如果 $c_{i,k}=0,d_{i,k}=5$，就是放电 5 kW。这样写的好处是每个变量都有明确范围，求解器更容易处理。

\subsection{静态电池约束}

静态电池能量更新为：
\begin{equation}
E^b_{i,k+1}
=
E^b_{i,k}
+\eta_b c_{i,k}\Delta t
-\frac{d_{i,k}\Delta t}{\eta_b}.
\end{equation}

逐项解释：
\begin{itemize}
  \item $E^b_{i,k}$：当前电池里有多少能量；
  \item $\eta_b c_{i,k}\Delta t$：这一小步充进去的有效能量；
  \item $\frac{d_{i,k}\Delta t}{\eta_b}$：这一小步放电导致电池内部减少的能量；
  \item $\eta_b<1$ 表示充放电有损耗。
\end{itemize}

SoC 和能量的关系为：
\begin{equation}
E^b_{i,k}=SOC^b_{i,k}C^b_i.
\end{equation}

例如容量 $C^b_i=10$ kWh，SoC 为 50\%，那么电池能量就是 5 kWh。

电池能量不能超过上下界：
\begin{equation}
SOC^b_{\min}C^b_i
\le
E^b_{i,k}
\le
SOC^b_{\max}C^b_i.
\end{equation}

充放电功率也不能超过最大值：
\begin{equation}
0\le c_{i,k}\le P^{b,\max}_i,
\qquad
0\le d_{i,k}\le P^{b,\max}_i.
\end{equation}

这些约束的物理意义是：电池不能充爆，不能放空，也不能瞬间以超过硬件能力的功率充放电。

\subsection{PV 在 MPC 中如何处理}

PV 表示光伏发电。对 MPC 来说，未来 PV 不是控制器能决定的，而是预测输入：
\begin{equation}
G_{i,t+k}.
\end{equation}

也就是说，MPC 不能命令太阳多发电或少发电。它只能决定：如果预测有这么多 PV，哪些用掉，哪些丢掉。丢掉的部分叫弃光，变量是：
\begin{equation}
u_{i,k}.
\end{equation}

弃光范围为：
\begin{equation}
0\le u_{i,k}\le \max(G_{i,t+k},0).
\end{equation}

实际使用的 PV 功率是：
\begin{equation}
G^{use}_{i,k}=G_{i,t+k}-u_{i,k}.
\end{equation}

如果 $u_{i,k}=0$，表示 PV 全部使用。如果 $u_{i,k}=2$ kW，表示有 2 kW 光伏没有进入负荷、电池或电网，而是被主动削减。弃光本身不是好事，但在某些情况下它能帮助满足功率平衡或避免不希望的出口。

\subsection{EV 约束}

EV 充电功率为 $e_{i,k}$。如果 EV 在家：
\begin{equation}
0\le e_{i,k}\le P^{ev,\max}_i.
\end{equation}

如果 EV 不在家：
\begin{equation}
e_{i,k}=0.
\end{equation}

这不是算法偏好，而是现实条件：车不在家时，家用充电桩无法给车充电。

EV 能量更新为：
\begin{equation}
E^{ev}_{i,k+1}
=
E^{ev}_{i,k}
+\eta_{ev}e_{i,k}\Delta t.
\end{equation}

EV 能量和 SoC 的关系为：
\begin{equation}
E^{ev}_{i,k}=SOC^{ev}_{i,k}C^{ev}_i.
\end{equation}

EV 能量也有上下界：
\begin{equation}
SOC^{ev}_{\min}C^{ev}_i
\le
E^{ev}_{i,k}
\le
SOC^{ev}_{\max}C^{ev}_i.
\end{equation}

\subsection{局部功率平衡}

对一个用户来说，调度后的净负荷可以写成：
\begin{equation}
n_{i,k}
=
L_{i,t+k}
-G_{i,t+k}
+u_{i,k}
+c_{i,k}
-d_{i,k}
+e_{i,k}.
\end{equation}

这个公式很重要，可以逐项读：
\begin{itemize}
  \item $L$：家庭负荷，要用电，所以增加净负荷；
  \item $-G$：PV 发电可以抵消负荷，所以减少净负荷；
  \item $+u$：弃掉 PV 后，可用光伏变少，所以净负荷增加；
  \item $+c$：电池充电要吸收电力，所以增加净负荷；
  \item $-d$：电池放电供给家庭，所以减少净负荷；
  \item $+e$：EV 充电要用电，所以增加净负荷。
\end{itemize}

如果 $n_{i,k}>0$，用户从电网买电。如果 $n_{i,k}<0$，用户向电网送电。

代码中还会用进口变量 $P^{imp}_{i,k}$ 和出口变量 $P^{exp}_{i,k}$ 表示：
\begin{equation}
P^{imp}_{i,k}-P^{exp}_{i,k}=n_{i,k}.
\end{equation}

同时会用二进制变量避免同一时间既进口又出口。直觉上就是：一个用户同一瞬间不应该既买电又卖电，否则数学上可能出现非物理的套利解。

\subsection{Local MPC 的目标函数}

Local MPC 的目标是最小化未来窗口内的总成本：
\begin{equation}
\min J_i^{local}.
\end{equation}

本项目中的主要形式可以写成：
\begin{equation}
J_i^{local}
=
\sum_{k=0}^{H-1}
\left[
\tilde p_{t+k}\Delta t\,c_{i,k}
-\tilde p_{t+k}\Delta t\,d_{i,k}
+\tilde p_{t+k}\Delta t\,e_{i,k}
+\epsilon_b\Delta t(c_{i,k}+d_{i,k})
\right]
+w^{ev}_{dep}g_i^2.
\end{equation}

逐项解释：
\begin{itemize}
  \item $\tilde p\Delta t\,c$：电池充电需要买电，是成本；
  \item $-\tilde p\Delta t\,d$：电池放电可以抵消高价购电，所以在成本里是负号；
  \item $\tilde p\Delta t\,e$：EV 充电也要买电，是成本；
  \item $\epsilon_b(c+d)$：很小的电池吞吐惩罚，用来避免无意义的频繁充放电；
  \item $w^{ev}_{dep}g_i^2$：soft 模式下 EV 离家 SoC 不够时的惩罚。
\end{itemize}

这里的 $g_i$ 是 EV 离家 SoC 缺口。它越大，说明车离家时越没充够。

\subsection{Local MPC 的优点和局限}

Local MPC 的优点：
\begin{itemize}
  \item 比 rule-based 更会利用未来电价、负荷和 PV 预测；
  \item 每个 agent 的优化问题比较小，求解相对快；
  \item 行为容易解释，因为它的目标函数和约束都很明确。
\end{itemize}

Local MPC 的局限：
\begin{itemize}
  \item 每个 agent 只顾自己，不知道其他 agent 同时在做什么；
  \item 多个用户可能同时在低价时充电，造成变压器压力；
  \item 它不直接把完整电压、线路电流等 AC 电网约束写进优化问题。
\end{itemize}

\section{ADMM MPC：让多个用户边自私边协调}

\subsection{为什么需要 ADMM}

假设一个小区里有 3 户家庭。每户都运行 Local MPC。某个时刻电价很低，于是 3 户都决定给电池和 EV 充电。对每户来说，这个决定都很合理；但对整个小区变压器来说，总功率可能太大。

这就是 Local MPC 的典型问题：
\begin{quote}
局部最优的动作，加在一起不一定是系统安全的动作。
\end{quote}

ADMM MPC 的目的就是在不把所有用户完全合成一个巨大优化问题的情况下，让它们对系统级约束进行协调。本项目里主要协调的是变压器聚合功率。

\subsection{先用故事理解 ADMM}

ADMM 可以先不看公式，理解成三轮对话：
\begin{enumerate}
  \item 每个家庭先各自做计划：我想在未来每一步用多少电；
  \item 协调器检查总和：所有家庭加起来有没有超过变压器限制；
  \item 如果超过了，就给每个家庭一个修正信号，让下一轮计划更接近安全总功率。
\end{enumerate}

这个过程会重复多轮。每一轮家庭仍然关心自己的电费，但也会受到协调信号影响。最后得到的计划既考虑局部经济性，也尽量满足系统总功率约束。

\subsection{ADMM 中三个新概念}

ADMM MPC 会出现三个对初学者陌生的量：
\begin{itemize}
  \item local variable：每个 agent 自己算出来的净负荷计划；
  \item consensus variable：协调后的公共副本，本文记为 $z_{i,k}$；
  \item dual variable：历史修正信号，本文记为 $u^{admm}_{i,k}$。
\end{itemize}

其中 $z_{i,k}$ 不是 PV 弃光，PV 弃光仍然是 $u_{i,k}$。为了避免混淆，本文件把 ADMM 的对偶变量写成 $u^{admm}_{i,k}$。

可以把 $u^{admm}$ 理解为“上几轮你偏离协调目标的累计提醒”。如果某个 agent 总是让总功率变得太高，这个提醒会在后续优化中推动它改变计划。

\subsection{被协调的量：净负荷贡献}

ADMM MPC 先定义每个 agent 的净负荷：
\begin{equation}
n_{i,k}
=
L_{i,t+k}-G_{i,t+k}+u_{i,k}+c_{i,k}-d_{i,k}+e_{i,k}.
\end{equation}

当前实现中的贡献系数为：
\begin{equation}
\alpha_{i,k}=1.
\end{equation}

所以每个 agent 对变压器总功率的贡献就是：
\begin{equation}
q_{i,k}=\alpha_{i,k}n_{i,k}=n_{i,k}.
\end{equation}

ADMM MPC 希望每个未来步都满足：
\begin{equation}
-\bar P^{trafo}
\le
\sum_{i=1}^{N} q_{i,k}
\le
\bar P^{trafo}.
\end{equation}

其中：
\begin{equation}
\bar P^{trafo}=0.97P^{trafo}_{limit}.
\end{equation}

这里的 0.97 是安全裕度。它表示不要把变压器刚好用到 100\%，而是提前留一点余量。

\subsection{局部子问题：每个 agent 怎样被协调}

在 ADMM MPC 中，每个 agent 的优化目标从 Local MPC 的 $J_i^{local}$ 变成：
\begin{equation}
\min
J_i^{local}
+
\frac{\rho}{2}
\left\|
\bm{\alpha}_i\odot\bm{n}_i-\bm{z}_i+\bm{u}^{admm}_i
\right\|_2^2.
\end{equation}

这个公式初看很吓人，但可以拆开：
\begin{itemize}
  \item $J_i^{local}$：自己原来的电费和 EV 惩罚；
  \item $\bm{\alpha}_i\odot\bm{n}_i$：自己计划对变压器的贡献；
  \item $\bm{z}_i$：协调器希望它接近的贡献副本；
  \item $\bm{u}^{admm}_i$：历史偏差带来的修正；
  \item $\rho$：协调项有多强。
\end{itemize}

平方项的物理意义是：如果 agent 的计划和协调目标差得太远，它就会被惩罚。$\rho$ 越大，agent 越愿意为了系统协调牺牲自己的局部电费最优。

\subsection{聚合投影：把总功率拉回安全范围}

每一轮局部求解后，协调器拿到所有 agent 的贡献 $q_{i,k}$。它要保证：
\begin{equation}
-\bar P^{trafo}
\le
\sum_i z_{i,k}
\le
\bar P^{trafo}.
\end{equation}

所有满足这个条件的 $\bm{z}$ 构成一个安全集合：
\begin{equation}
\mathcal{C}
=
\left\{
\bm{z}:
-\bar P^{trafo}\le \sum_i z_{i,k}\le \bar P^{trafo},
\ \forall k
\right\}.
\end{equation}

投影可以写成：
\begin{equation}
\bm{z}^{new}
=
\Pi_{\mathcal{C}}(\bm{q}+\bm{u}^{admm}).
\end{equation}

这里的“投影”可以理解成“找离原计划最近的安全计划”。如果总功率没有超过限制，投影前后基本不变。如果超过了限制，投影会把超出的部分分配回各个 agent 的副本，使总和落回安全区间。

例如某一步 3 个 agent 的贡献加起来是 110 kW，但安全上限是 100 kW，就超了 10 kW。项目实现中的思想是把这 10 kW 的超出量按简单方式分摊回各个 agent 的协调副本。这样下一轮每个 agent 再优化时，就会感受到“总功率太高，需要往下调”的信号。

\subsection{对偶变量更新：记住偏差}

投影之后，ADMM 更新对偶变量：
\begin{equation}
\bm{u}^{admm}
\leftarrow
\bm{u}^{admm}
+
\bm{q}
-
\bm{z}^{new}.
\end{equation}

如果 $\bm{q}$ 和 $\bm{z}^{new}$ 相差很大，说明 agent 自己想要的计划和系统安全计划差距大。这个差距会被记进 $u^{admm}$，影响下一轮局部优化。

所以 ADMM 不是一次命令所有人服从，而是通过多轮“计划、检查、修正、再计划”逐渐达成一致。

\subsection{什么时候停止迭代}

ADMM 用两个残差判断是否已经足够一致：
\begin{equation}
r^{pri}
=
\|\bm{q}-\bm{z}\|_2,
\end{equation}
\begin{equation}
r^{dual}
=
\rho\|\bm{z}^{new}-\bm{z}^{old}\|_2.
\end{equation}

原始残差 $r^{pri}$ 衡量 agent 自己的计划和协调计划是否接近。对偶残差 $r^{dual}$ 衡量协调变量本身是否还在大幅变化。

当二者都小到一定程度时：
\begin{equation}
r^{pri}\le \epsilon^{pri},
\qquad
r^{dual}\le \epsilon^{dual},
\end{equation}
算法就认为已经收敛。收敛不是说“宇宙真理已经找到”，而是说继续迭代带来的改变已经很小，可以停止了。

代码中还会做 residual balancing，也就是根据两个残差的相对大小调整 $\rho$。直觉是：如果局部计划和协调计划差得太多，就加强协调；如果协调变量自己抖得太厉害，就减弱一点。

\subsection{第一步动作投影：为什么还要再修一次}

ADMM 算的是未来 $H$ 步，但环境马上只执行第 1 步。因此项目代码在 ADMM 迭代后，还会对第一步动作做一个小的二次规划投影。

它的目标是：
\begin{equation}
\min
\sum_i
\left[
(P^b_i-\hat P^b_i)^2
+(u_i-\hat u_i)^2
+(e_i-\hat e_i)^2
\right].
\end{equation}

其中带帽子的 $\hat P^b_i,\hat u_i,\hat e_i$ 是 ADMM 原本给出的动作。这个小问题的意思是：
\begin{quote}
尽量不要改太多，但必须把当前要执行的动作修到主要安全约束内。
\end{quote}

特别是第一步变压器聚合功率要满足：
\begin{equation}
-\bar P^{trafo}
\le
\sum_i n_{i,0}
\le
\bar P^{trafo}.
\end{equation}

这一步很重要，因为即使未来预测步还没有完全协调好，当前马上执行的动作也应该尽量满足主要聚合约束。

\subsection{ADMM MPC 的优点和局限}

ADMM MPC 的优点：
\begin{itemize}
  \item 比 Local MPC 多了变压器聚合功率协调；
  \item 不需要把所有 agent 完全合成一个巨大优化问题；
  \item 可以减少多个 agent 同时充电造成的系统峰值风险。
\end{itemize}

ADMM MPC 的局限：
\begin{itemize}
  \item 当前主要协调变压器聚合功率，不等于完整 AC 潮流安全约束；
  \item 电压和线路约束没有像 Global MISOCP 那样完整进入优化；
  \item 有限迭代内可能还没完全收敛；
  \item 为了系统安全，可能牺牲一部分局部经济性。
\end{itemize}

\section{soft 和 hard 模式：到底软在哪里，硬在哪里}

\subsection{先给核心直觉}

soft 和 hard 主要讨论 EV 离家 SoC 要求。假设车早上要出门，希望离家时至少有 80\% 电量。

soft 模式的态度是：
\begin{quote}
你最好达到 80\%，如果没达到，就在目标函数里罚你。
\end{quote}

hard 模式的态度是：
\begin{quote}
只要物理上还来得及，就必须保证离家目标；如果控制器给的充电太少，执行层会把它抬高。
\end{quote}

因此：
\begin{itemize}
  \item soft 是“允许违反，但违反要付代价”；
  \item hard 是“尽量不允许违反，必要时直接修正动作”。
\end{itemize}

\subsection{soft 模式的数学表达}

设离家步在当前预测窗口中的位置为 $k_{dep}$. EV 离家 SoC 缺口定义为：
\begin{equation}
g_i
\ge
SOC^{ev}_{req}
-
\frac{E^{ev}_{i,k_{dep}}}{C^{ev}_i},
\qquad
g_i\ge 0.
\end{equation}

也可以写成更直观的形式：
\begin{equation}
g_i
=
\max\left(
0,
SOC^{ev}_{req}
-
SOC^{ev}_{i,k_{dep}}
\right).
\end{equation}

如果车离家时已经达到要求，$g_i=0$。如果只充到 70\%，而要求是 80\%，那么 $g_i=10\%$。

soft 模式把缺口加入目标函数：
\begin{equation}
J
\leftarrow
J
+
w^{ev}_{dep}g_i^2.
\end{equation}

平方 $g_i^2$ 的意义是：小缺口会被罚，大缺口会被罚得更重。权重 $w^{ev}_{dep}$ 越大，优化器越不愿意让 EV 离家电量不足。

soft 模式的好处是灵活，不容易因为要求太严格而导致优化问题不可行。坏处是它没有绝对保证：如果惩罚权重不够大，或者预测窗口没有看到离家时刻，EV 仍可能没充够。

\subsection{hard 模式在本项目 MPC 中的实现}

在标准优化教材里，hard 模式常常会直接写成硬约束：
\begin{equation}
E^{ev}_{i,k_{dep}}
\ge
SOC^{ev}_{req}C^{ev}_i.
\end{equation}

这句话表示：离家时 EV 能量必须达到目标，不能通过“付罚金”逃避。

但本项目的 Local MPC 和 ADMM MPC 要特别区分：它们的局部优化问题本身仍然主要使用 soft 缺口结构。进入 hard 实验配置时，\texttt{scripts/mpc.py} 会启用环境执行层的 EV hard projection：
\begin{equation}
\texttt{ev\_hard\_projection\_enabled=True}.
\end{equation}

也就是说，对 Local MPC 和 ADMM MPC 来说，hard 模式的关键不是优化目标里多了一个神奇公式，而是环境在执行动作前会检查 EV 是否再不充就来不及。如果来不及，就把 EV 充电功率强制抬高到最低必要值。

\subsection{hard projection 怎么算}

设当前 EV SoC 为 $SOC^{ev}_{i,t}$，目标离家 SoC 为 $SOC^{ev}_{req}$。环境会估计离家前还剩多少个可充电时间步，记为：
\begin{equation}
K_t.
\end{equation}

如果未来每一步都用最大功率充电，未来最多还能增加：
\begin{equation}
\Delta SOC^{ev,\max}_{i,t}
=
\frac{K_tP^{ev,\max}_i\Delta t\eta_{ev}}{C^{ev}_i}.
\end{equation}

如果现在 SoC 加上未来最多可补的 SoC 仍然不够，那么当前这一步就必须充一点。当前最低必要充电功率为：
\begin{equation}
P^{ev,min}_{i,t}
=
\max\left(
0,
\frac{
\left(
SOC^{ev}_{req}
-SOC^{ev}_{i,t}
-\Delta SOC^{ev,\max}_{i,t}
\right)C^{ev}_i
}
{\Delta t\eta_{ev}}
\right).
\end{equation}

控制器原本请求的 EV 充电功率是 $P^{ev,raw}_{i,t}$。hard projection 后实际执行：
\begin{equation}
P^{ev,exec}_{i,t}
=
\operatorname{clip}
\left(
\max(P^{ev,raw}_{i,t},P^{ev,min}_{i,t}),
0,
P^{ev,\max}_i
\right).
\end{equation}

这条公式可以翻译成三句话：
\begin{itemize}
  \item 如果控制器请求已经足够大，就按控制器的来；
  \item 如果控制器请求太小，而再不充就来不及，环境会抬高到 $P^{ev,min}$；
  \item 最后还要限制在 $0$ 到 $P^{ev,\max}$ 之间，不能超过物理最大功率。
\end{itemize}

\subsection{soft 和 hard 行为差异}

soft 模式下，MPC 会在两件事之间权衡：
\begin{equation}
\text{现在少充电省钱}
\quad \text{vs.}\quad
\text{离家 SoC 不足被罚}.
\end{equation}

hard 模式下，控制器可以等待便宜电价，但只能等到“不再等也来得及”的边界：
\begin{equation}
\text{可以等待低价}
\quad \text{直到}\quad
\text{再等就无法按时达标}.
\end{equation}

所以 hard 模式通常会带来：
\begin{itemize}
  \item EV 离家 SoC 满足率更稳定；
  \item 对惩罚权重不那么敏感；
  \item 接近离家时可能出现强制补电；
  \item 如果低价窗口已经错过，强制补电可能发生在高价时段；
  \item EV 强制充电可能和电池套利、变压器约束产生冲突。
\end{itemize}

\section{Local MPC 和 ADMM MPC 在 soft/hard 下的对比}

\begin{longtable}{p{3.1cm}p{5.7cm}p{5.7cm}}
\toprule
项目 & Local MPC & ADMM MPC \\
\midrule
\endhead
基本思想 & 每个 agent 独立省钱 & 每个 agent 省钱，同时通过 ADMM 协调总功率 \\
未来信息 & 使用未来价格、负荷、PV 预测 & 使用未来价格、负荷、PV 预测 \\
电池约束 & 有 SoC 和功率约束 & 有 SoC 和功率约束 \\
PV 处理 & PV 是预测输入，可通过弃光变量削减 & PV 是预测输入，可通过弃光变量削减 \\
EV soft & 离家缺口进入目标函数惩罚 & 离家缺口进入局部目标函数惩罚 \\
EV hard & 主要依赖环境 hard projection 修正执行动作 & 主要依赖环境 hard projection 修正执行动作 \\
变压器约束 & 不显式协调多 agent 总功率 & 显式协调聚合变压器功率 \\
电压/线路约束 & 不完整显式建模 & 不完整显式建模，主要处理聚合功率 \\
主要优点 & 简单、快、容易解释 & 更能处理多用户同时动作导致的聚合风险 \\
主要风险 & 各自最优可能叠加成系统过载 & 协调可能未完全收敛，也可能牺牲部分经济性 \\
\bottomrule
\end{longtable}

\section{Global MISOCP 和两种 MPC 的关系}

仓库中的 \texttt{MISOCP} 是一个更强的全局优化基线。它不像 Local MPC 那样每个用户分开算，也不像 ADMM MPC 那样只协调聚合变压器功率，而是把更完整的电网约束放进同一个优化模型。

它包含类似 DistFlow 的网络约束，例如支路 $(u,v)$ 上：
\begin{equation}
P_{\ell,k}
=
\sum_{m\in\mathcal{C}(v)}P_{m,k}
+r_{\ell}\ell_{\ell,k}
+p^{inj}_{v,k},
\end{equation}
\begin{equation}
Q_{\ell,k}
=
\sum_{m\in\mathcal{C}(v)}Q_{m,k}
+x_{\ell}\ell_{\ell,k},
\end{equation}
\begin{equation}
v_{v,k}
=
v_{u,k}
-2(r_{\ell}P_{\ell,k}+x_{\ell}Q_{\ell,k})
+(r_{\ell}^2+x_{\ell}^2)\ell_{\ell,k}.
\end{equation}

二阶锥松弛可以写成：
\begin{equation}
P_{\ell,k}^2+Q_{\ell,k}^2
\le
v_{u,k}\ell_{\ell,k}.
\end{equation}

对初学者来说，不需要马上掌握 MISOCP 的全部数学。只要先记住：
\begin{itemize}
  \item Local MPC：每户自己优化；
  \item ADMM MPC：每户自己优化，但协调总变压器功率；
  \item Global MISOCP：中央调度器直接考虑更完整的电网物理约束。
\end{itemize}

\section{从零基础角度总结}

理解 MPC，可以抓住四个问题：
\begin{enumerate}
  \item 控制器能决定什么？答案是电池充放电、EV 充电、PV 弃光；
  \item 控制器不能违反什么？答案是电池容量、充放电功率、EV 连接状态、SoC 上下界等物理约束；
  \item 控制器想优化什么？答案是未来电费、EV 离家缺口惩罚、少量电池动作惩罚；
  \item 控制器有没有考虑别人？Local MPC 没有，ADMM MPC 通过变压器聚合功率考虑了一部分。
\end{enumerate}

因此，本项目中的 MPC 可以理解为：
\begin{quote}
每一步，控制器拿着未来预测数据，先在纸上安排未来几步怎么充电、放电和弃光；然后只执行当前一步；下一步再根据新状态重新安排。
\end{quote}

soft 和 hard 的区别可以理解为：
\begin{quote}
soft 是“没充够就罚分”；hard 是“快来不及时直接把充电动作抬高”。
\end{quote}

带着这个直觉再回头看公式，公式就不再是一堆符号，而是在描述真实系统里的能量流、费用和安全边界。

\end{document}
